problem:  
Let \((M^{2n}, g, J)\) be a compact Kähler manifold with **nonnegative quadratic bisectional curvature** (QBC), meaning  
\[
\sum_{i,j} R_{i\bar{i} j \bar{j}}(x_i - x_j)^2 \ge 0
\]
for every unitary frame \(\{e_i\}\) and every real vector \((x_1,\ldots,x_n)\), where \(R_{i\bar{i} j\bar{j}}\) denotes the curvature components with respect to this frame.  
It is known that positive QBC lies strictly between nonnegative orthogonal bisectional curvature and nonnegative Ricci curvature, but its geometric and topological implications remain poorly understood.  

**Open problem:**  
Assume \(M\) has nonnegative QBC and positive QBC at least at one point. If \(X^p, Y^q \subset M\) are closed complex submanifolds satisfying \(p+q \ge n\), must \(X \cap Y \neq \varnothing\)?  
If not necessarily, determine the precise curvature-threshold condition—possibly involving positivity of QBC in rank \( \ge k\)—that still guarantees the Frankel-type intersection property.  

Why is it a "good" problem:  
This problem replaces the artificial averaging notions with the **intrinsically natural curvature invariant (quadratic bisectional curvature)**, which interpolates geometrically between Ricci and full bisectional curvature while preserving essential directional data. It asks whether Frankel’s deep intersection phenomenon persists under this intermediate curvature positivity, thus probing the minimal curvature conditions enforcing topological intersection rigidity in Kähler manifolds. Resolving it would clarify the boundary between curvature positivity sufficient to control geodesic behavior and that needed to enforce intersection properties, potentially leading to a new curvature–topology principle parallel to Frankel’s theorem. The formulation is concise yet touches on uncharted terrain linking analytic curvature inequalities, complex geometry, and global topology—a hallmark of a “good” research problem.